\input{tools/book_preamble.tex}
\newcommand{\BookTitle}{操作系统}
\hypersetup{pdftitle={操作系统·408自学讲义},pdfsubject={概念、方法与代表题}}
\begin{document}
\pagenumbering{alph}
\begin{titlepage}
\thispagestyle{empty}
\vspace*{25mm}
{\sffamily\color{accent}\fontsize{16}{24}\selectfont 408 自学讲义}\par
\vspace{15mm}
{\sffamily\bfseries\color{ink}\fontsize{32}{45}\selectfont 操作系统}\par
\vspace{15mm}
\noindent\rule{\textwidth}{0.5pt}\par
\vspace{5mm}
\noindent 例题就近解析，自测题在章末给出答案。\par
\vspace{5mm}\noindent 正文内容审阅稿：本轮尚未进行版式检查。\par
\vfill
\noindent A4 单栏 · 2026 年 10 月\par
\end{titlepage}
\pagenumbering{roman}
\tableofcontents
\clearpage
\phantomsection\chapter*{阅读约定}
\addcontentsline{toc}{chapter}{阅读约定}
例题中的【年份】表示 408 真题；〖节号-题号〗表示王道书相应节的习题编号。
改动原题条件时标为“改编”；补充题另标“原创”。先根据题目条件辨认模型，再阅读连续的计算或执行过程。

本讲义默认下标、指针、位数和计量单位在初次使用处定义。
同一个字母在不同主题中可能表示不同的量，以当前定义为准。
并发过程、地址转换和通信时序中的结论都依赖给定假设；题目改变条件时，应重新检查解法成立的前提。

正文末尾附题目索引与方法索引，可从目录或书签跳转；索引右侧为正文页码。
\clearpage
\pagenumbering{arabic}
\setcounter{chapter}{0}
\input{build/revision3/OS/body.tex}
\input{build/revision3/OS/indexes.tex}
\end{document}
